\chapter{CONCLUSION}
An effective simple EEG based BCI system was developed for typing based on blink and focus signals. Thus, it was proven that a simple BCI system may be designed with even limited amount of channels and simple equipment. Filter design was implemented and tested for isolation of the frequency range containing useful signals, amplification, filtering and FFT analysis were performed to investigate the EEG signals. The overall performance of the system was investigated using confusion matrix and class performance measures.

At the same time, there are a few limitations identified during the testing which influence the system performance. They include the usage of the same channel for blink and focus detection, the usage of a matrix board instead of a PCB, the presence of a fixed gain at DRL stage, limitations of Arduino processing power and distortions caused by non-ideal filter.

Thus, the current project demonstrates the practical implementation of biomedical instrumentation including biopotential signal acquisition, analog front-end design, active filters, noise reduction and suppression, grounding and signal conditioning and digital signal processing techniques. In conclusion, it can be stated that the current project demonstrates practical application of brain signal measurements.